\section{My role}

\textbf{1a. Roles allocated vs taken.} At the wk12 kickoff I was explicitly allocated the computer-vision workstream: train the dummy detector, own \texttt{vision.py}, and feed detections to the mission script Robin was writing. That was the whole of my formal remit. By wk22 I had also, without anyone assigning it, written the SITL simulator, the state machine in \texttt{main.py}, the progressive flight-test ladder (\texttt{tests/flight/0a}--\texttt{4}), the geofence, the web ground station, Demetro's FDR speaking scripts, and the draft of the team's complaint letter to the course organiser. The drift from ``CV lead'' to ``everything infrastructure'' is the single most important fact about my project, and I want to be honest about why it happened.

It was not ambition. Around wk14, after the first flight-day cancellation, I noticed the team was becoming idle waiting for hardware, and my reflex was to build a way to keep working. I now see this is because my instinct under uncertainty is to convert anxiety into code --- if I am writing a simulator I am not thinking about the fact that we have no drone. By wk17 the simulator was running full missions and I had started refactoring \texttt{main.py}; by wk20, after the third cancelled flight day (session 2026-03-30), I was writing unit tests and a modularity audit of my own architecture. Each individual step felt like duty. The aggregate shape was a teammate quietly absorbing four roles because doing so felt safer than asking why the team's coordination was not catching the gap. I initially framed this as ``the team needs infrastructure and I can build it'' --- in fact it was closer to ``I trust my own hands more than I trust a conversation I have not had,'' and I now see I preferred the first framing because it protected me from noticing the second.

\textbf{1b. Most proud of.} The contribution I am most proud of is the simulation framework, specifically the decision at wk14 to build it at all. Technically it is \texttt{simple\_simulator.py} (2508 lines), the SITL integration, and the dual-backend \texttt{vision.py} that lets the same mission code run on a laptop with Ultralytics and on a Pi with TFLite without a single \texttt{if} branch in the state machine. The reason I am proud of it is not the line count. It is that in mid-March we had three cancelled flight days and zero seconds of real flight data, and the simulator was the only reason the project produced any verifiable engineering output at the D6 milestone. Without it the eight weeks between wk14 and wk22 would have been dead time; with it we validated the full state machine, retrained the detector against synthetic data, and walked the R01--R12 requirements one by one until all twelve were met in simulation (session 2026-04-03).

What I am actually proud of is that I made a judgement call about risk under ambiguity. In wk14 nobody was asking for a simulator --- the official plan still assumed the drone would fly the following week. Building it meant committing 40--60 hours to something that would be unnecessary if the hardware worked. The uncomfortable realisation, which I could not articulate at the time, was that I no longer trusted the plan. I now recognise that the simulator was partly a technical insurance policy and partly a psychological one: it gave the team a way to make progress visible to ourselves when the official project had stalled. I had previously thought of engineering judgement as picking the right technique. This project forced me to reframe that --- judgement is also choosing when to stop trusting the stated constraints and build a second path, and being willing to own the time cost if the first path recovers.

\textbf{1c. Focus differently in future.} The honest lesson is not ``I should work less.'' It is that I should work less \emph{alone}. I grew up inside a Ukrainian engineering culture where the default unit of delivery is one person who owns the problem end-to-end, and the heroic pattern is the engineer who stays at the bench until the thing runs. That instinct served me on this project but it also let me quietly take work off teammates who would have grown from doing it themselves --- I wrote Demetro's speaking scripts four times instead of sitting with him while he wrote them once, and I refactored \texttt{main.py} from 1411 to 754 lines in a weekend rather than pairing with Robin on the parts he would later integrate.

In my next team project I will cap personal solo-work at roughly 60\% of my hours and invest the other 40\% in explicit pairing and teaching sessions, and the way I will know whether it is working is a specific falsifiable signal: whether a teammate can independently complete a handoff task I would previously have done myself, within 48 hours of our pairing session, without messaging me for help. That is a measurable test, not a wish. The harder commitment underneath it is that I will treat the discomfort of slow collaboration as the actual skill I am trying to build --- not as overhead on top of the real work. I leave this project believing that engineering leadership is less about being the person who can deliver any module alone and more about being the person who makes sure the module does not need to be delivered alone. The moment this became clear to me was session 2026-04-03, when I ran the modularity audit on my own codebase and scored \texttt{main.py} at 3.0/5.0 --- the weakest module in the project was the one I had spent the most hours on by myself.
